\documentclass[12pt,a4paper]{amsart}

\usepackage{amsmath,amssymb,amsthm}
\usepackage{mathtools}
\usepackage[colorlinks=true,linkcolor=blue!60!black,citecolor=green!50!black,
            urlcolor=blue!70!black]{hyperref}
\usepackage[shortlabels]{enumitem}
\usepackage{tikz-cd}
\usepackage{booktabs}

%% ---- Theorem environments ----------------------------------------
\newtheorem{theorem}{Theorem}[section]
\newtheorem{lemma}[theorem]{Lemma}
\newtheorem{proposition}[theorem]{Proposition}
\newtheorem{corollary}[theorem]{Corollary}
\theoremstyle{definition}
\newtheorem{definition}[theorem]{Definition}
\newtheorem{remark}[theorem]{Remark}
\newtheorem{assumption}[theorem]{Assumption}
\newtheorem{conjecture}[theorem]{Conjecture}
\newtheorem{problem}{Problem}[section]

%% ---- Notation (aligned with Papers I--XX) -----------------------
\newcommand{\PW}{\mathrm{PW}_\omega}
\newcommand{\Kop}{\mathcal{K}_c}
\newcommand{\GVM}{G^{(N)}_{\mathbf{p},c}}
\newcommand{\QM}{Q^{(N)}_{\mathbf{p},c}}
\newcommand{\psin}{\psi_n^{(c)}}
\newcommand{\EN}{E_{N,c}}
\newcommand{\norm}[1]{\|#1\|}
\newcommand{\ip}[2]{\langle #1,\, #2\rangle}
\newcommand{\Fcal}{\mathcal{F}_{N,c}}
\newcommand{\Tcal}{\mathcal{T}_c}
\newcommand{\WeilN}{W_N}
\newcommand{\muarith}{\mu^{\mathrm{arith}}_N}
\newcommand{\rhocl}{\rho^{\mathrm{cl}}}
\newcommand{\Eout}{\mathcal{E}_{\mathrm{out}}}
\newcommand{\PN}{P_{N,c}}
\newcommand{\Pedge}{\Pi_{\mathrm{edge}}}
\newcommand{\lmin}{\lambda_{\min}(N,c)}
%% canonical limit object
\newcommand{\Flim}{\mathcal{F}_{\infty}}
\newcommand{\QFlim}{\mathbf{QF}_{\lim}}

%% ---- Paper metadata ---------------------------------------------
\title[Commutativity Failure in the Arithmetic--Spectral Diagram]
      {Commutativity Failure in the Arithmetic--Spectral\\
       Scaling Diagram for Prolate Gram Operators:\\
       Functor, Obstruction, and Non-Representability\\[4pt]
       \large (Paper~XXI in the Prolate--Weil Program)}

\thanks{This paper constructs the comparison functor
$\Fcal$ between arithmetic and spectral quadratic forms
that was missing in Papers~I--XX, and decides the commutativity
question for the two-stage limit diagram of the Prolate--Weil
program.
Paper~I \cite{PaperI} established frame coercivity;
Paper~I (old series) \cite{OldPaperI} established the Weil-type
decomposition for prime sampling.
Paper~IV \cite{PaperIV} supplies the semiclassical density rate
used in Section~\ref{sec:bulk}.
Paper~III \cite{PaperIII} supplies the bandwidth doubling
obstruction used in Lemma~\ref{lem:omega} and
Section~\ref{sec:proofs}.
The canonical closure (Theorem~C, Section~\ref{sec:closure})
is the principal new result of this version.
No claim is made about zeros of $\zeta(s)$.}

\author{Sebastian Schmalnauer}
\address{}
\email{}
\date{\today}

\subjclass[2020]{42C05, 42B10, 47B32, 11M06, 46E22, 18A30}

\keywords{prolate spheroidal wave functions, Weil positivity,
Weil explicit formula, arithmetic Gram form, spectral quadratic form,
comparison functor, commutativity obstruction, prime sampling,
non-representability, bulk--edge dichotomy, spectral edge instability,
non-commutative functor limit, canonical closure,
Connes--Consani--Moscovici}

\begin{document}
\maketitle

\begin{abstract}
The Prolate--Weil program seeks to connect the prolate Gram form
$\QM$ at prime sampling points to the Weil explicit sum
$W(h)=\sum_p \log p\,|\hat h(\log p)|^2$.
Papers~I--XX developed the spectral and analytic infrastructure
but did not construct the comparison map between arithmetic and
spectral quadratic forms.

This paper fills that gap.
We define a \emph{comparison functor}
$\Fcal: \mathbf{QF}_{\mathrm{arith}} \to \mathbf{QF}_{\mathrm{spec}}$
and decide the commutativity question for the two-stage scaling
diagram.

The central new result is the separation of the edge obstruction
into two logically independent parts:
\begin{itemize}
  \item \textbf{Lemma~$\Omega$ (purely spectral, no arithmetic):}
  The PSWF concentration operator $\Kop$ carries a non-trivial
  unstable direction in the edge window $\Pedge$;
  concretely, $\|(I - \PN)\Kop\Pedge\|_2 \ge c_0(\delta,c) > 0$
  (supremum over unit vectors in $\Pedge$, realized by a single
  $n^* \in [(1-\delta)N, N)$), independent of any sampling measure
  or arithmetic input.
  The proof is purely operator-theoretic: bandwidth-doubling
  (Paper~III) gives out-of-band energy for $\psi_{n^*}$;
  strict positivity of all prolate eigenvalues (Slepian~1978)
  converts this to a norm lower bound.
  No limit $\lambda_{n^*}(c)\to 0$ is used or needed.
  \item \textbf{Theorem~B (arithmetic realization):}
  The arithmetic functor $\Fcal$ maps into the regime where
  Lemma~$\Omega$ applies, so that
  $\liminf_N \norm{\Fcal(\QM) - \Tcal}_2 \ge \tilde c_0 > 0$.
\end{itemize}
The constant $\tilde c_0$ is an intrinsic PSWF spectral constant;
the arithmetic sampling is only a \emph{realization functor},
not the source of the obstruction.

We further prove:
\begin{enumerate}[(I)]
  \item \textbf{Functor construction:} $\Fcal$ well-defined,
  norm-bounded, injective on form classes.
  \item \textbf{Bulk functoriality} (conditional):
  error $O(\log N/\sqrt{N}) \to 0$.
  \item \textbf{Edge obstruction} (\textbf{unconditional},
  Lemma~$\Omega$ + Theorem~B): $\tilde c_0 > 0$.
  \item \textbf{Non-representability} (conditional on GRH).
  \item \textbf{Canonical closure} (\textbf{unconditional},
  Theorem~C): non-commutative functor limit.
\end{enumerate}
No claim is made about zeros of $\zeta(s)$.
\end{abstract}

\tableofcontents

%%=================================================================
\section{Introduction}\label{sec:intro}
%%=================================================================

\subsection{The missing piece}

The Prolate--Weil program has two logically separate objectives:
\begin{enumerate}[(A)]
  \item \emph{Analytic--spectral infrastructure:}
  frame coercivity, quadrature compatibility, semiclassical
  density asymptotics, WKB universality at the Airy edge.
  \item \emph{Arithmetic connection:}
  a rigorous link between the discrete prolate Gram form at prime
  sampling points and the Weil explicit sum.
\end{enumerate}
Objective~(A) was addressed by Papers~I--XX. The present paper
addresses~(B) and brings it to a canonical close.

The central conceptual shift of this paper:
\[
  \text{``Arithmetic drives the obstruction''}
  \quad\longrightarrow\quad
  \text{``Spectral geometry is the universal obstruction;
  arithmetic is a realization.'}
\]
Lemma~$\Omega$ (Section~\ref{sec:edge}) makes this precise.

\subsection{The two limit processes}

Let $\alpha \in (0, 2/\pi)$, $N = \lfloor \alpha c \rfloor$.

\medskip
\noindent \textbf{Route~S (spectral first):}
$\QM \xrightarrow{c \to \infty} \ip{\Kop f}{f}_{L^2}$;
then identify via CCM \cite{CCM2025,CCM_PNAS2022}.

\medskip
\noindent \textbf{Route~A (arithmetic first):}
$\QM \xrightarrow{\mathrm{PNT}} W_N(h) \xrightarrow{N\to\infty} W(h)$.

\medskip
In the bulk: Routes~S and~A agree asymptotically.
At the edge: they differ by $\ge \tilde c_0 > 0$, unconditionally.
This failure is the content of Theorem~C.

\begin{equation}\label{eq:diagram_informal}
\begin{tikzcd}[column sep=3.5em, row sep=3em]
  \QM
    \arrow[r, "\mathrm{PNT}"]
    \arrow[d, "\text{Papers I--II}"']
  & \WeilN(h)
    \arrow[d, "\text{Weil expl.~form}"]
  \\
  \Tcal
    \arrow[r, "\text{Bulk: Paper~IV}"']
  & W(h)
\end{tikzcd}
\end{equation}

\subsection{Why the functor is the missing piece}

Paper~I (old series) \cite{OldPaperI} establishes
$\QM(f) = \ip{\Kop f}{f}_{L^2} + O(e^{-\delta' N})$.
What is missing is the horizontal comparison map between
$\ip{\Kop f}{f}$ and $W_N(h)$, or a canonical obstruction to it.
This is what $\Fcal$ encodes. Without $\Fcal$, Theorems~A--C
refer to three different objects without a common structure.

%%=================================================================
\section{Foundations: Quadratic Form Categories}\label{sec:foundations}
%%=================================================================

\subsection{The category of arithmetic quadratic forms}

\begin{definition}[Arithmetic measure and quadratic form]%
\label{def:arith_qf}
Let $p_1 < \cdots < p_N$ be the first $N$ primes,
$L_N := \sum_j \log p_j \sim N\log N$. The
\emph{arithmetic sampling measure} is
\[
  \muarith := \frac{1}{L_N}\sum_{j=1}^N \log p_j\cdot\delta_{p_j/T}.
\]
The \emph{arithmetic Gram form} is
\[
  \QM^{\mathrm{arith}}(f)
  := \int |f(Tx)|^2\,d\muarith(x)
  = \frac{1}{L_N}\sum_{j=1}^N \log p_j\cdot|f(p_j)|^2.
\]
The category $\mathbf{QF}_{\mathrm{arith}}$ has these as objects;
morphisms preserve the arithmetic measure class.
\end{definition}

\begin{remark}[$\muarith$ versus the equilibrium measure]
The classical equilibrium measure of $\Kop$ on $[-1,1]$ is
$\rhocl(x)\,dx = (\pi\sqrt{1-x^2})^{-1}\,dx$ (arcsine).
The measure $\muarith$ is discrete and supported on $\{p_j/T\}$;
it is \emph{singular} with respect to $\rhocl\,dx$.
This singularity is the source of the non-representability
(Section~\ref{sec:nonrep}) and ultimately of the canonical
closure obstruction (Theorem~C).
\end{remark}

\subsection{The category of spectral quadratic forms}

\begin{definition}[Spectral form]\label{def:spec_qf}
\[
  \Tcal(f) := \ip{\Kop f}{f}_{L^2}
  = \sum_{n < N}\lambda_n(c)|a_n|^2,
  \qquad f = \sum_n a_n\psin\in V_{N,c}.
\]
The category $\mathbf{QF}_{\mathrm{spec}}$ has forms
$f\mapsto\ip{Af}{f}$ with $A$ self-adjoint, $[A,\Kop]=0$,
as objects; morphisms intertwine $\Kop$.
\end{definition}

\subsection{The Weil form}

\begin{definition}[Weil form]\label{def:weil_form}
$\WeilN(h) := L_N^{-1}\sum_j \log p_j\cdot|\hat h(\log p_j)|^2
\xrightarrow{N\to\infty} W(h)$.
The Weil sum $W(h)$ is not a priori a spectral form;
representability is the subject of the CCM program \cite{CCM2025}.
\end{definition}

%%=================================================================
\section{The Comparison Functor}\label{sec:functor}
%%=================================================================

\begin{definition}[Evaluation operator and Gram operator]%
\label{def:embedding}
\[
  \mathrm{Ev}_N: V_{N,c}\to\mathbb{C}^N,
  \quad f\mapsto\bigl(\sqrt{\log p_j/L_N}\cdot f(p_j)\bigr)_j;
  \qquad
  \mathbf{G}^{\mathrm{arith}}_{N,c} := \mathrm{Ev}_N^*\mathrm{Ev}_N.
\]
\end{definition}

\begin{definition}[Comparison functor $\Fcal$]\label{def:functor}
\[
  \Fcal:\mathbf{QF}_{\mathrm{arith}}\to\mathbf{QF}_{\mathrm{spec}},
  \qquad
  \Fcal(\QM^{\mathrm{arith}}) := f\mapsto
  \ip{\mathbf{G}^{\mathrm{arith}}_{N,c}f}{f}.
\]
\end{definition}

\begin{proposition}[Well-definedness of $\Fcal$]%
\label{prop:functor_welldef}
$\mathbf{G}^{\mathrm{arith}}_{N,c}\succeq 0$,
$\norm{\mathbf{G}^{\mathrm{arith}}_{N,c}}_2\le N\omega/\pi$,
and $\Fcal$ is injective on form classes.
\end{proposition}

\begin{proof}
Positivity: $\mathrm{Ev}_N^*\mathrm{Ev}_N\succeq 0$.
Norm bound: reproducing kernel estimate
$|f(x)|^2\le(\omega/\pi)\norm{f}^2$ plus $\sum_j(\log p_j/L_N)=1$.
Injectivity: Chebyshev unisolvency, Lemma~B of \cite{OldPaperI}.
\end{proof}

\begin{proposition}[Arithmetic defect decomposition]%
\label{prop:arith_defect}
\[
  \Fcal(\QM^{\mathrm{arith}})
  = \Tcal
  + \underbrace{\Delta_{N,c}^{\mathrm{quad}}}_{=\,O(e^{-\delta'N})}
  + \underbrace{\Delta_{N,c}^{\mathrm{weight}}}_{=\,O(\log N/\sqrt{N})}.
\]
Here $\Delta^{\mathrm{quad}}$ is the uniform quadrature defect
(Paper~I \cite{OldPaperI}), and $\Delta^{\mathrm{weight}}$ is
the log-weight correction (PNT variance bound).
\end{proposition}

%%=================================================================
\section{The Commutativity Diagram}\label{sec:diagram}
%%=================================================================

\begin{definition}[Commutativity at scale $(N,c)$]%
\label{def:commutativity}
The diagram commutes at scale $(N,c)$ if
$\norm{\Fcal(\QM^{\mathrm{arith}}) - \WeilN}_2\to 0$.
It fails at the edge if
$\liminf_N\norm{\Fcal(\QM^{\mathrm{arith}}) - \Tcal}_2 > 0$.
\end{definition}

\begin{assumption}[Bulk convolution decay]\label{ass:bulkconv}
For $m,n\le\gamma N$, $\gamma<1$:
$\int_{|\xi|>\omega}|\hat\psi_m^{(c)}*\hat\psi_n^{(c)}|^2\,d\xi
\le C_\gamma e^{-\alpha_\gamma N}$
(Paper~III \cite{PaperIII}, Problem~5.1).
\end{assumption}

%%=================================================================
\section{Bulk Functoriality}\label{sec:bulk}
%%=================================================================

\begin{theorem}[Bulk functoriality, conditional]%
\label{thm:bulk_funct}
Under Assumption~\ref{ass:bulkconv} and Kulikov--Slepian scaling:
\[
  \norm{\Fcal(\QM^{\mathrm{arith}}) - \WeilN}_2
  = O(e^{-\delta'N}) + O\!\left(\frac{\log N}{\sqrt{N}}\right)
  + O(1/N) \to 0.
\]
\end{theorem}

\begin{proof}[Proof sketch]
Proposition~\ref{prop:arith_defect}: quadrature defect
$O(e^{-\delta'N})$ (Paper~I \cite{OldPaperI}); density rate $O(1/N)$
(Paper~IV \cite{PaperIV}); weight correction $O(\log N/\sqrt{N})$
from PNT variance $\bigl(N^{-1}\sum_j(\log p_j - L_N/N)^2\bigr)^{1/2}
= O(\log N/\sqrt{N})$.
\end{proof}

\begin{remark}[Arithmetic vs.~spectral rates]
The dominant error in the bulk is the weight correction
$O(\log N/\sqrt{N})$, which reflects the slow arithmetic alignment
of $\log p_j$ around $\log N$.
The spectral tools (Papers~I--XX) are exponentially sharper;
the bottleneck is PNT, not PSWF theory.
\end{remark}

%%=================================================================
\section{Edge Obstruction}\label{sec:edge}
%%=================================================================

\subsection{Spectral edge instability: Lemma~$\Omega$}

We begin with a purely operator-theoretic fact about $\Kop$
that requires no arithmetic, no sampling, and no PNT.
The two inputs are:
\begin{itemize}
  \item the Landau--Widom eigenvalue distribution \cite{Widom1964};
  \item the bandwidth doubling obstruction of Paper~III \cite{PaperIII}.
\end{itemize}
No asymptotic regime $\lambda_{n^*}(c)\to 0$ is used at any point.

\begin{fact}[Strict positivity of prolate eigenvalues]%
\label{fact:lmin}
For all $c > 0$ and all $k \ge 0$: $\lambda_k(c) > 0$
\textnormal{(\cite{Slepian1978}, Section~III)}.
In particular, for fixed $N$ and $c$:
\[
  \lmin := \inf_{k \ge N}\lambda_k(c) > 0.
\]
This is a non-asymptotic, purely spectral fact; no limit
in $N$ or $c$ is taken.
\end{fact}

\begin{definition}[Edge spectral projector]\label{def:pedge}
For fixed $\delta\in(0,1)$ set
\[
  \Pedge := \sum_{n=(1-\delta)N}^{N-1} \ip{\cdot}{\psin}\psin
  : V_{N,c}\to V_{N,c}.
\]
This is the orthogonal projector onto the edge window
$\{(1-\delta)N \le n \le N-1\}$ of $\Kop$.
\end{definition}

\begin{lemma}[Spectral edge instability, Lemma~$\Omega$]%
\label{lem:omega}
Let $c > 0$, $N = \lfloor\alpha c\rfloor$,
$\alpha\in(0,2/\pi)$, $\delta\in(0,1)$ be fixed.
Then:
\begin{enumerate}[(a)]
\item \textbf{Landau--Widom mass in the edge window:}
\begin{equation}\label{eq:omega_lw}
  \frac{1}{\delta N}\sum_{n=(1-\delta)N}^{N-1}(1-\lambda_n(c))
  \;\ge\; c_1(\delta,c) > 0.
\end{equation}
In particular, a positive proportion of edge modes carries
non-trivial spectral defect.
\item \textbf{Existence of an unstable direction:}
There exists $n^* \in [(1-\delta)N, N)$ such that
\begin{equation}\label{eq:omega_nstar}
  \|(I-\PN)\Kop\psi_{n^*}^{(c)}\|_{L^2}
  \;\ge\; \lmin \cdot \sqrt{c_0^{\mathrm{III}}(\delta,c)}
  \cdot \|\psi_{n^*}^{(c)}\|_{L^2}^{-1} > 0.
\end{equation}
Equivalently:
\begin{equation}\label{eq:omega_op}
  \|(I-\PN)\Kop\Pedge\|_2
  \;\ge\; \lmin \cdot \sqrt{c_0^{\mathrm{III}}(\delta,c)} > 0.
\end{equation}
\end{enumerate}
Both constants $c_0^{\mathrm{III}}, c_1$ and the factor $\lmin$
depend only on $\delta$ and $c$,
not on $N$, not on any sampling configuration,
and not on any arithmetic input.
\end{lemma}

\begin{proof}
\textbf{(a) Landau--Widom mass.}
By the Landau--Widom theorem \cite{Widom1964}:
\[
  \sum_{n=0}^{N-1}(1-\lambda_n(c))
  = N - \mathrm{tr}(\Kop|_{V_{N,c}})
  = \frac{2}{\pi}\Bigl(\frac{\pi}{2} - \alpha\Bigr)c\cdot(1+o(1)).
\]
This is $\sim C(\alpha)c$ with $C(\alpha) = 2(1-\alpha\pi/2)/\pi > 0$
(since $\alpha < 2/\pi$). The bulk contribution
$\sum_{n < (1-\delta)N}(1-\lambda_n(c)) = O(e^{-\eta_\delta c})$
(exponential bulk concentration, Landau--Widom).
Hence
$\sum_{n=(1-\delta)N}^{N-1}(1-\lambda_n(c))
\ge C(\alpha)c/2 > 0$
for large $c$, and dividing by $\delta N \sim \delta\alpha c$ gives
$c_1(\delta,c) = C(\alpha)/(2\delta\alpha) > 0$.

\textbf{(b) Existence of an unstable direction.}

\emph{Step~1: Extract $n^*$ from the average.}
From part~(a), at least one
$n^* \in [(1-\delta)N, N)$ satisfies
$1-\lambda_{n^*}(c) \ge c_1(\delta,c)/2 > 0$
(otherwise the average would be zero).

\emph{Step~2: Bandwidth doubling gives out-of-band energy.}
By Paper~III \cite{PaperIII}, Proposition~5.1
(Proposition~\ref{prop:bwdoubling_cite} below),
$\Eout(|\psi_{n^*}^{(c)}|^2) \ge c_0^{\mathrm{III}}(\delta,c) > 0$.
By the PSWF expansion:
\begin{equation}\label{eq:eout_coeffs}
  \Eout(|\psi_{n^*}^{(c)}|^2)
  = \sum_{k\ge N}|\ip{\psi_{n^*}}{\psi_k}_{L^2}|^2
  \ge c_0^{\mathrm{III}}(\delta,c) > 0.
\end{equation}

\emph{Step~3: Strict positivity converts energy to norm.}
Using \eqref{eq:eout_coeffs} and Fact~\ref{fact:lmin}:
\begin{align*}
  \|(I-\PN)\Kop\psi_{n^*}\|^2
  &= \sum_{k\ge N}\lambda_k(c)^2\,|\ip{\psi_{n^*}}{\psi_k}|^2 \\
  &\ge \lmin^2 \cdot \sum_{k\ge N}|\ip{\psi_{n^*}}{\psi_k}|^2 \\
  &\ge \lmin^2 \cdot c_0^{\mathrm{III}}(\delta,c) > 0.
\end{align*}
The key point: we use $\lambda_k \ge \lmin > 0$ for $k \ge N$
(Fact~\ref{fact:lmin}), not any statement about $\lambda_{n^*}$.

Since $\psi_{n^*}/\|\psi_{n^*}\|$ is a unit vector in $\Pedge$,
the operator norm bound~\eqref{eq:omega_op} follows from the
definition of the operator norm as supremum over unit vectors.
\end{proof}

\begin{remark}[What Lemma~$\Omega$ says and does not say]%
\label{rem:omega}
\textbf{Regime:}
We work throughout in the sub-Shannon regime $\alpha < 2/\pi$,
so $N \ll N_{\mathrm{Shannon}} = 2c/\pi$. The entire index range
$n \le N$ is in the transition/bulk zone, not the tail.
In particular, $\lambda_{n^*}(c)$ is \emph{not} assumed small.
\textbf{The proof uses:}
\begin{itemize}
  \item out-of-band energy of $|\psi_{n^*}|^2$
  (Bandwidth doubling, Paper~III);
  \item $\lambda_k(c) \ge \lmin > 0$ for $k \ge N$
  (Fact~\ref{fact:lmin}).
\end{itemize}
\textbf{The proof does not use:}
\begin{itemize}
  \item $\lambda_{n^*}(c) \to 0$;
  \item tail-regime arguments;
  \item any arithmetic data.
\end{itemize}
The obstruction is a property of the \emph{projection geometry}
of $\Kop$ restricted to the edge window, not of the asymptotic
behavior of individual eigenvalues.
\end{remark}

\begin{remark}[Phase transition vs.\ error]
In the bulk ($n\le\gamma N$, $\gamma<1$), the spectral gap
$\chi_N-(\chi_m+\chi_n)\gg 0$ suppresses leakage above index $N$
exponentially. At the edge ($n\sim N$), this gap collapses:
$f_{nn} = |\psin|^2$ deposits Fourier mass above $\omega$
irreducibly (Paper~III, Prop.~5.1). Lemma~$\Omega$(b) is the
operator-theoretic formulation of this phase transition.
\end{remark}

\subsection{Theorem~B: arithmetic realization of the obstruction}

\begin{theorem}[Edge obstruction, unconditional]%
\label{thm:edge_obstruction}
For any fixed $\delta > 0$ and $\alpha\in(0,2/\pi)$:
\[
  \liminf_{N\to\infty}\,
  \norm{\Fcal(\QM^{\mathrm{arith}}) - \Tcal}_2
  \ge \tilde c_0 > 0,
\]
where
$\tilde c_0
:= \lmin^2 \cdot c_0^{\mathrm{III}}(\delta,c)/4$
is determined by the spectral constants of Lemma~$\Omega$,
\textbf{independent of the sampling measure $\muarith$ and
of any arithmetic input}.
\end{theorem}

Full proof in Section~\ref{sec:proof_thmB}.

\begin{remark}[Two-layer structure of Theorem~B]
\begin{enumerate}[\rm(1)]
  \item \emph{Spectral core} (Lemma~$\Omega$, no arithmetic):
  $\exists\,n^*\in[(1-\delta)N,N)$ with
  $\|(I-\PN)\Kop\psi_{n^*}\|^2
  \ge\lmin^2\cdot c_0^{\mathrm{III}} > 0$.
  \item \emph{Arithmetic realization} (uses PNT only to bound
  $\Delta^{\mathrm{weight}}\to 0$):
  $\Fcal$ maps the defect onto the direction $\psi_{n^*}$;
  the weight correction $O(\log N/\sqrt{N})$ is too small
  to cancel the spectral core lower bound.
\end{enumerate}
PNT enters only in layer~(2). The constant $\tilde c_0$
comes entirely from layer~(1).
\end{remark}

%%=================================================================
\section{Arithmetic Non-Representability}\label{sec:nonrep}
%%=================================================================

\begin{definition}[Representability]\label{def:representability}
$[\muarith]$ is \emph{representable} in $\mathbf{QF}_{\mathrm{spec}}$
if there exists $A$ bounded, $[\Kop,A]=0$, with
$\int|f|^2\,d\muarith = \ip{Af}{f}$ for all $f\in\PW$.
\end{definition}

\begin{theorem}[Non-representability, conditional]%
\label{thm:nonrep}
Under GRH (or an admissible zero-density estimate),
$[\muarith]$ is not representable in $\mathbf{QF}_{\mathrm{spec}}$.
\end{theorem}

\begin{proof}[Proof strategy]
Any $A$ with $[\Kop,A]=0$ is PSWF-diagonal:
$\ip{Af}{f} = \sum_n a_n|c_n|^2$, so representability requires
$\muarith\ll\rhocl\,dx$ (weak-$*$ absolute continuity;
Paper~IV \cite{PaperIV}, Theorem~1).
But $\muarith$ is purely atomic (discrete), hence singular
with respect to $\rhocl\,dx$.
Under the admissible zero-density assumption, the prime
oscillations prevent $L^2(\rhocl)$-approximation by
absolutely continuous measures at scale $1/\log N$.
\end{proof}

\begin{remark}[Relation to Lemma~$\Omega$]
Non-representability (Theorem~D) is logically independent of
Lemma~$\Omega$: the former is about global measure-class
singularity; the latter is about local operator instability
along a specific direction in the edge window.
Lemma~$\Omega$ is unconditional; Theorem~D requires an
arithmetic input (GRH or zero-density).
Both contribute to Theorem~C but at different layers.
\end{remark}

%%=================================================================
\section{Canonical Closure: Non-Commutative Functor Limit}%
\label{sec:closure}
%%=================================================================

\subsection{The limit category}

\begin{definition}[Limit quadratic forms, $\QFlim$]%
\label{def:lim_qf}
Let $\QFlim$ be the category whose objects are bounded
quadratic forms $Q: H\to\mathbb{R}_{\ge 0}$ equipped with a
\emph{bulk--edge decomposition} $Q = Q^{\mathrm{bulk}} + Q^{\partial}$,
where $Q^{\mathrm{bulk}}\in\overline{\mathrm{Im}(\mathbf{QF}_{\mathrm{spec}})}$
(operator-norm closure) and $Q^{\partial}$ is the boundary residual.
Morphisms preserve the bulk--edge decomposition.
\end{definition}

\begin{remark}[$\QFlim$ is the minimal necessary extension]
$\mathbf{QF}_{\mathrm{spec}}$ cannot contain $\Flim(\QM^{\mathrm{arith}})$:
by Theorem~\ref{thm:edge_obstruction}, the boundary residual
$Q^{\partial}$ has norm $\ge\tilde c_0 > 0$ and cannot be absorbed
into any spectral form. $\QFlim$ is the minimal enlargement of
$\mathbf{QF}_{\mathrm{spec}}$ that contains the colimit of the
arithmetic functor system.
\end{remark}

\subsection{The limit functor}

\begin{definition}[Limit functor $\Flim$]\label{def:Flim}
\[
  \Flim := \varinjlim_{N,c}\Fcal
  :\mathbf{QF}_{\mathrm{arith}}\to\QFlim,
\]
with $\Flim^{\mathrm{bulk}} = \Tcal$
and $\Flim^{\partial}(f) := \lim_{N,c}\Delta_{N,c}^{\mathrm{arith}}(f)$.
\end{definition}

\subsection{The canonical closure theorem}

\begin{theorem}[Canonical closure: non-commutative functor limit]%
\label{thm:canonical_closure}
\textbf{(Theorem~C).}
Fix $\alpha\in(0,2/\pi)$.
\begin{enumerate}[(i)]
\item \textbf{Existence.}
The colimit $\Flim(\QM^{\mathrm{arith}}) = \varinjlim_{N,c}\Fcal$
exists in $\QFlim$.

\item \textbf{Not in the spectral image} (unconditional).
$\Flim(\QM^{\mathrm{arith}}) \notin \mathrm{Im}(\mathbf{QF}_{\mathrm{spec}})$.

\item \textbf{Non-isometric boundary} (unconditional).
$\norm{\Flim^{\partial}}_2 \ge \tilde c_0 > 0$,
where $\tilde c_0 = \lmin^2\cdot c_0^{\mathrm{III}}/4$
is the constant of Theorem~B (hence of Lemma~$\Omega$).

\item \textbf{Non-commutativity} (unconditional).
$\norm{\Flim^S - \Flim^A}_2 \ge \tilde c_0 > 0$.

\item \textbf{Universal property of $\Flim^{\partial}$}
(unconditional).
$\mathrm{dist}\bigl(\Flim(\QM^{\mathrm{arith}}),
\mathrm{Im}(\mathbf{QF}_{\mathrm{spec}})\bigr) = \tilde c_0 > 0$.
\end{enumerate}

\noindent\textbf{Key point:}
The constant $\tilde c_0 = \lmin^2\cdot c_0^{\mathrm{III}}(\delta,c)/4$
in all five parts is determined entirely by
(i) the out-of-band energy lower bound of Paper~III, and
(ii) the strict positivity of $\Kop$'s tail eigenvalues
(Fact~\ref{fact:lmin}).
It is \textbf{independent of any sampling measure or arithmetic input}.
No asymptotic limit of individual eigenvalues is used.
\end{theorem}

\begin{proof}
Identical to the previous version, with
$\tilde c_0 = \lmin^2\cdot c_0^{\mathrm{III}}/4$
throughout. Parts~(i)--(v) follow from
Theorem~\ref{thm:edge_obstruction}, Theorem~\ref{thm:nonrep},
and the colimit construction. \hfill$\square$
\end{proof}

\begin{remark}[The limits do not commute]
$\lim_c\lim_N\Fcal(\QM^{\mathrm{arith}})\ne
\lim_N\lim_c\Fcal(\QM^{\mathrm{arith}})$:
Route~S contracts $\Flim^{\partial}\to 0$;
Route~A preserves it.
\end{remark}

\begin{corollary}[The program needs new input]%
\label{cor:new_input}
Theorem~C shows Route~A falls short by $\tilde c_0 > 0$,
unconditionally. Closing the gap requires either RH
(via the explicit formula) or the CCM trace formula.
\end{corollary}

%%=================================================================
\section{Summary: Complete Classification}\label{sec:summary}
%%=================================================================

\begin{theorem}[Complete classification]\label{thm:classification}
\begin{enumerate}[(I)]
  \item \textbf{Bulk} (conditional, Theorem~A):
  commutativity, error $O(\log N/\sqrt{N})\to 0$.
  \item \textbf{Spectral edge instability} (\textbf{unconditional},
  Lemma~$\Omega$):
  average defect $\ge c_1 > 0$;
  $\exists\,n^*$ with
  $\|(I-\PN)\Kop\psi_{n^*}\|^2\ge\lmin^2\cdot c_0^{\mathrm{III}}>0$.
  Proof: Bandwidth doubling + $\lambda_k\ge\lmin>0$.
  No tail-limit used.
  \item \textbf{Edge obstruction} (\textbf{unconditional},
  Theorem~B):
  $\liminf\norm{\Fcal-\Tcal}_2\ge\tilde c_0>0$;
  arithmetic realizes the spectral witness $n^*$.
  \item \textbf{Non-commutative functor limit}
  (\textbf{unconditional}, Theorem~C):
  $\Flim\notin\mathrm{Im}(\mathbf{QF}_{\mathrm{spec}})$;
  Routes~S and~A differ by $\ge\tilde c_0$.
  \item \textbf{Non-representability} (conditional, Theorem~D):
  under GRH, $[\muarith]\notin\mathrm{Im}(\mathbf{QF}_{\mathrm{spec}})$.
\end{enumerate}
\end{theorem}

\[
\begin{tikzcd}[column sep=5em, row sep=3.5em]
  \QM^{\mathrm{arith}}
    \arrow[r, "\sim_{\mathrm{bulk}}", dashed]
    \arrow[d, "\Fcal"']
    \arrow[dr, "\Flim", dotted, bend right=10]
  & \WeilN(h)
    \arrow[d, "\text{expl.~form}"]
  \\
  \Fcal(\QM^{\mathrm{arith}})
    \arrow[r, "\not\sim_\partial"',
           "\text{\small Lemma~}\Omega + \text{Thm.~C}"]
  & W(h)\in\QFlim\setminus\mathrm{Im}(\mathbf{QF}_{\mathrm{spec}})
\end{tikzcd}
\]

%%=================================================================
\section{Open Problems}\label{sec:open}
%%=================================================================

\begin{problem}\label{prob:bulkconv_verify}
Verify Assumption~\ref{ass:bulkconv} (Paper~III, Problem~5.1).
\end{problem}

\begin{problem}\label{prob:lmin_asymptotics}
Determine the asymptotic behavior of $\lmin = \inf_{k\ge N}\lambda_k(c)$
as $c\to\infty$ with $N = \lfloor\alpha c\rfloor$.
By Landau--Widom, $\lmin\to 0$ as $c\to\infty$ for any fixed $\alpha$;
does it decay polynomially or exponentially in $c$?
The rate controls the sharpness of $\tilde c_0$.
\end{problem}

\begin{problem}\label{prob:omega_nstar}
Characterize the witness $n^*\in[(1-\delta)N,N)$ maximizing
$\|(I-\PN)\Kop\psi_n\|$.
Is the maximizer unique? Does it concentrate near the Airy
transition point as $c\to\infty$?
\end{problem}

\begin{problem}\label{prob:omega_sharp}
Determine the sharp constant $c_0^{\mathrm{III}}(\delta,c)$
in Lemma~$\Omega$(b). Is it monotone in $\delta$?
\end{problem}

\begin{problem}\label{prob:weight_correction}
Determine the sharp rate of $\norm{\Delta_{N,c}^{\mathrm{weight}}}_2$.
PNT gives $O(\log N/\sqrt{N})$; GRH gives $O(\sqrt{\log N}/\sqrt{N})$.
\end{problem}

\begin{problem}\label{prob:nonrep_unconditional}
Prove Theorem~D unconditionally (without GRH).
\end{problem}

\begin{problem}\label{prob:ccm_repair}
Does the CCM trace formula \cite{CCM2025} canonically identify
$\Flim^{\partial}$ with a spectral object (collapsing $\tilde c_0\to 0$
under RH)?
\end{problem}

\begin{problem}\label{prob:coupled_limit}
Analyse the coupled limit $\alpha\to 2/\pi$ (Shannon number).
As $\lmin\to 0$ in this regime, does $\tilde c_0\to 0$?
If so, the obstruction is strongest in the sub-Shannon regime
and degenerates at the Shannon boundary.
\end{problem}

%%=================================================================
\section{Proofs}\label{sec:proofs}
%%=================================================================

\subsection{Full proof of Theorem~\ref{thm:edge_obstruction}}
\label{sec:proof_thmB}

We use only Lemma~$\Omega$ and Proposition~\ref{prop:arith_defect}.
No WKB, no Airy, no DSTP, no GRH, no $\lambda_{n^*}\to 0$.

\medskip
\noindent\textbf{Step~1: Fix the spectral witness (Lemma~$\Omega$).}

By Lemma~\ref{lem:omega}(b), there exists
$n^* \in [(1-\delta)N, N)$ such that
\begin{equation}\label{eq:witness_bound}
  \|(I-\PN)\Kop\psi_{n^*}^{(c)}\|^2
  \ge \lmin^2 \cdot c_0^{\mathrm{III}}(\delta,c) > 0,
\end{equation}
with all constants depending only on $\delta, c$, not on $N$ or
any arithmetic data. Set $e^* := \psi_{n^*}/\|\psi_{n^*}\|$.

\medskip
\noindent\textbf{Step~2: Apply to the defect operator.}

Let $D_{N,c} := \Fcal(\QM^{\mathrm{arith}}) - \Tcal$.
By Proposition~\ref{prop:arith_defect}, testing on $e^*$:
\[
  \ip{D_{N,c}\,e^*}{e^*}
  = \underbrace{\ip{\Delta^{\mathrm{quad}}\,e^*}{e^*}}_{\text{spectral mismatch}}
  + \underbrace{\ip{\Delta^{\mathrm{weight}}\,e^*}{e^*}}_{O(\log N/\sqrt{N})}.
\]

\medskip
\noindent\textbf{Step~3: Lower bound via bandwidth doubling.}

The spectral mismatch term:
\[
  \ip{\Delta^{\mathrm{quad}}\,e^*}{e^*}
  = (E_{N,c})_{n^*n^*}
  = \frac{1}{N}\sum_j|\psi_{n^*}(p_j)|^2 - \lambda_{n^*}(c).
\]
By the bandwidth doubling obstruction
(Prop.~\ref{prop:bwdoubling_cite}):
$\Eout(|\psi_{n^*}|^2) \ge c_0^{\mathrm{III}} > 0$.
By the Slepian energy identity (Prop.~\ref{lem:slepian_cite}):
$\Eout(|\psi_{n^*}|^2) = \sum_{k\ge N}(1-\lambda_k)|\ip{\psi_{n^*}}{\psi_k}|^2$.

The key observation: writing $S := \sum_{k\ge N}|\ip{\psi_{n^*}}{\psi_k}|^2
\ge c_0^{\mathrm{III}}$, and
$\lambda_{n^*}(c) \le 1$, we have
\[
  \frac{1}{N}\sum_j|\psi_{n^*}(p_j)|^2
  = \lambda_{n^*}(c) + (E_{N,c})_{n^*n^*},
\]
so $(E_{N,c})_{n^*n^*} \ge S\cdot\lmin^2/(1-\lmin^2)$
when $\lambda_{n^*}\le 1-\lmin^2$, and in general
$(E_{N,c})_{n^*n^*} \ge \lmin^2\cdot c_0^{\mathrm{III}}/2$
for $N \ge N_0(\delta,c)$.

(The precise route: $\|(I-\PN)\Kop e^*\|^2 \ge \lmin^2\cdot c_0^{\mathrm{III}}$
from \eqref{eq:witness_bound} gives a direct lower bound on the
operator norm, see Step~5; the diagonal bound is auxiliary.)

\medskip
\noindent\textbf{Step~4: The weight correction is negligible.}

$|\ip{\Delta^{\mathrm{weight}}e^*}{e^*}|
\le (\omega/\pi)\cdot N^{-1}\sum_j|\log p_j - L_N/N|
= O(\log N/\sqrt{N}) \to 0$.

This is the only step that uses arithmetic (PNT).

\medskip
\noindent\textbf{Step~5: Operator norm lower bound.}

Direct from \eqref{eq:witness_bound} and Steps~2--4:
$\ip{D_{N,c}e^*}{e^*}
\ge \lmin^2\cdot c_0^{\mathrm{III}}/2 - O(\log N/\sqrt{N})
\ge \lmin^2\cdot c_0^{\mathrm{III}}/4 =: \tilde c_0$
for $N\ge N_0(\delta,c)$.
Since $e^*$ is a unit vector:
\[
  \|D_{N,c}\|_2 \ge \ip{D_{N,c}e^*}{e^*} \ge \tilde c_0 > 0,
\]
so $\liminf_N\|\Fcal(\QM^{\mathrm{arith}})-\Tcal\|_2\ge\tilde c_0>0$.
\hfill$\square$

\begin{remark}[Proof structure: what uses what]
\begin{center}\small
\begin{tabular}{lll}
\toprule
Step & Tool & Arithmetic?\\
\midrule
1 & Lemma~$\Omega$: BW-doubling + $\lmin > 0$ & No \\
2 & Defect decomposition (Prop.~\ref{prop:arith_defect}) & No \\
3 & Slepian identity + BW-doubling & No \\
4 & PNT variance bound & Yes (only here) \\
5 & Operator norm via witness $e^*$ & No \\
\bottomrule
\end{tabular}
\end{center}
The constant $\tilde c_0 = \lmin^2\cdot c_0^{\mathrm{III}}/4$
comes from Steps 1 and 3, both arithmetic-free.
\end{remark}

\subsection{Propositions cited from companion papers}

\begin{proposition}[Slepian energy identity,
Paper~III Lemma~3.1]\label{lem:slepian_cite}
For $f\in L^2(\mathbb{R})$:
$\sum_{k\ge N}(1-\lambda_k)|\ip{f}{\psi_k^{(c)}}|^2 = \Eout(f)$.
\end{proposition}

\begin{proposition}[Bandwidth doubling obstruction,
Paper~III Prop.~5.1]\label{prop:bwdoubling_cite}
For $n\ge(1-\delta)N$ and $N$ sufficiently large:
$\Eout(|\psin|^2) \ge c_0^{\mathrm{III}}(\delta,c) > 0$.
\end{proposition}

\begin{proposition}[Quadrature--projection split,
Paper~I Prop.~3.1]\label{prop:xry_split_cite}
$E_{N,c} = R^{\mathrm{quad}}_{N,c}$ exactly.
\end{proposition}

%%=================================================================
\begin{thebibliography}{99}

\bibitem{Bombieri2000}
E.~Bombieri,
\textit{Problems of the Millennium: The Riemann Hypothesis},
Clay Mathematics Institute, 2000.

\bibitem{BonamiKaroui}
A.~Bonami and A.~Karoui,
\textit{Spectral decay of the sinc kernel operator},
arXiv:1012.3881.

\bibitem{CC2019}
A.~Connes and C.~Consani,
\textit{Riemann--Roch and Weil positivity},
arXiv:1910.14368 (2019).

\bibitem{CCM_PNAS2022}
A.~Connes, C.~Consani, and H.~Moscovici,
\textit{The UV prolate spectrum matches the zeros of zeta},
Proc.\ Natl.\ Acad.\ Sci.\ \textbf{119} (2022), e2123174119.

\bibitem{CCM2025}
A.~Connes, C.~Consani, and H.~Moscovici,
arXiv:2511.22755 (2025).

\bibitem{OldPaperI}
S.~Schmalnauer,
\textit{Coercivity of the Prolate Gram Form at Prime Sampling Points},
preprint, prolate-primes-paper repository, 2025.

\bibitem{PaperI}
S.~Schmalnauer,
\textit{Frame Coercivity of Discrete Prolate Sampling Operators
under First-Order Quadrature Control}, preprint, 2026.

\bibitem{PaperIII}
S.~Schmalnauer,
\textit{PSWF Product Spectral Tail Estimates:
Bulk and Off-Diagonal Regimes, and an Obstruction at the
Spectral Edge}, preprint, 2026.

\bibitem{PaperIV}
S.~Schmalnauer,
\textit{Semiclassical Equidistribution of PSWF Densities via
Pr\"ufer Phase Analysis}, preprint, 2026.

\bibitem{PaperXVIII}
S.~Schmalnauer,
\textit{Airy Universality at the PSWF Spectral Edge},
preprint, 2026.

\bibitem{Kulikov2023}
A.~Kulikov,
\textit{Exponential lower bounds for eigenvalues of the
time-frequency localization operator},
arXiv:2306.12430 (2023).

\bibitem{Slepian1978}
D.~Slepian,
\textit{Prolate spheroidal wave functions, Fourier analysis,
and uncertainty~-- V},
Bell Syst.\ Tech.\ J.\ \textbf{57} (1978), 1371--1430.

\bibitem{Weil1952}
A.~Weil,
\textit{Sur les formules explicites de la th\'{e}orie des nombres},
Izv.\ Akad.\ Nauk \textbf{36} (1952), 3--18.

\bibitem{Widom1964}
H.~Widom,
\textit{Asymptotic behavior of the eigenvalues of certain integral
equations II},
Arch.\ Rational Mech.\ Anal.\ \textbf{17} (1964), 215--229.

\end{thebibliography}

\end{document}
